% Source-aligned English for printed p.161.
\SourcePage{161}
\begingroup
\setlength{\abovedisplayskip}{0.5\baselineskip}
\setlength{\belowdisplayskip}{0.5\baselineskip}
\setlength{\abovedisplayshortskip}{0.35\baselineskip}
\setlength{\belowdisplayshortskip}{0.5\baselineskip}

\LayerTarget{tohoku:construction:divisor-sequence:p2}
\noindent $\mathbf D$ is flabby, since a section of $\mathbf D$ over a nonempty open
subset $U\subset X$ is a formal linear combination $\sum n_iV_i$ of
irreducible hypersurfaces $V_i$ in $U$, and is therefore the restriction of
the section $\sum n_iV_i$ of $\mathbf D$ on $X$. Consequently,

\LayerTarget{tohoku:eq:divisor-resolution}
\[
0\longrightarrow\mathbf O^*\longrightarrow\mathbf K^*\longrightarrow
\mathbf D\longrightarrow0\longrightarrow0\longrightarrow\cdots
\]
is a resolution of $\mathbf O^*$ by flabby sheaves ($\mathbf K^*$ is flabby
because it is constant and $X$ is irreducible; cf. the end of 3.3). We thus
obtain the values of $H^p(X,\mathbf O^*)$; the symbol $\Phi$ is omitted when
$\Phi$ is the family of all closed subsets of $X$:

\LayerTarget{tohoku:eq:divisor-class-cohomology}
\[
H^1(X,\mathbf O^*)=
\Gamma(\mathbf D)/\operatorname{Im}\bigl(\Gamma(\mathbf K^*)\bigr),
\]
the group of divisor classes modulo principal divisors. This familiar fact
also follows immediately from the exact cohomology sequence. Moreover,
$H^i(X,\mathbf O^*)=0$ for $i\geq2$, a result I first obtained in a much less
direct way by the method of Čech coverings. This result extends without
change when $X$ is a ``variety scheme'' in the sense of [5 bis], and more
generally to ``arithmetic varieties'' defined by gluing together ``spectra''
of commutative rings [8]. The following application can also be formulated
for arithmetic varieties.

\LayerTarget{tohoku:prop:3.4.1}
\textsc{Proposition 3.4.1.} --- \textit{Let $X$ be an irreducible algebraic
variety over an algebraically closed field $k$, whose local rings are
factorial; a nonsingular variety is an example. Then
$H^i(X,\mathbf O^*)=0$ for $i\geq2$. If $E$ is a locally trivial algebraic
fiber bundle over $X$ with structure group the projective group
$\operatorname{GP}(n-1,k)$ (cf. [20]), then $E$ is isomorphic to the fiber
space associated with a locally trivial algebraic fiber bundle whose
structure group is the linear group $\operatorname{Gl}(n,k)$.}

\LayerTarget{tohoku:proof:3.4.1}
For every algebraic group $G$, let $\mathbf O(G)$ denote the sheaf of groups
of germs of regular maps from $X$ to $G$. The first assertion of the
proposition is $H^i(X,\mathbf O(k^*))=0$ for $i\geq2$ and has already been
proved. In the terminology developed in [11], the second assertion says that
the canonical map

\LayerTarget{tohoku:eq:pgl-lift-map}
\[
H^1\bigl(X,\mathbf O(\operatorname{Gl}(n,k))\bigr)
\longrightarrow
H^1\bigl(X,\mathbf O(\operatorname{GP}(n-1,k))\bigr)
\]
is surjective. To prove this, consider the exact sequence of algebraic groups

\LayerTarget{tohoku:eq:central-extension}
\[
e\longrightarrow k^*\longrightarrow\operatorname{Gl}(n,k)
\longrightarrow\operatorname{GP}(n-1,k)\longrightarrow e,
\]
where the first homomorphism is the natural isomorphism of $k^*$ onto the
center of $\operatorname{Gl}(n,k)$. It is easy to check that the fibration of
$\operatorname{Gl}(n,k)$ by the subgroup $k^*$ is locally trivial, that is,
it has a rational section. The preceding exact sequence therefore gives an
exact sequence of sheaves

\LayerTarget{tohoku:eq:sheaf-central-extension}
\[
e\longrightarrow\mathbf O(k^*)
\longrightarrow\mathbf O(\operatorname{Gl}(n,k))
\longrightarrow\mathbf O(\operatorname{GP}(n-1,k))
\longrightarrow e,
\]
where $\mathbf O(k^*)$ lies in the center of
$\mathbf O(\operatorname{Gl}(n,k))$. The proposition now follows from
$H^2(X,\mathbf O(k^*))=0$ and from the corollary to the following result,
which generalizes [11, Proposition 5.7.2, Corollary], where we had to impose
paracompactness assumptions.

\LayerTarget{tohoku:prop:3.4.2:p1}
\textsc{Proposition 3.4.2.} --- \textit{Let $X$ be a topological space and
let $e\longrightarrow F\longrightarrow G\longrightarrow H\longrightarrow e$
be a}

\TohokuContinuation{162}
\endgroup
